%!TEX root = ../thesis.tex
%*******************************************************************************
%****************************** Chapter 6 **************************************
%*******************************************************************************

\chapter{CONCLUSION}
\label{chap:conclusion}

\ifpdf
    \graphicspath{{Chapter6/Figs/Raster/}{Chapter6/Figs/PDF/}{Chapter6/Figs/}}
\else
    \graphicspath{{Chapter6/Figs/Vector/}{Chapter6/Figs/}}
\fi

%********************************** Section 6.1 ********************************
\section{Summary}
\label{sec:con-summary}

This thesis studied how to adapt an open-source vision-language model to Vietnamese visual
question answering under tight resource constraints: parameter-efficient fine-tuning only,
and a single Colab GPU. The team chose Qwen2.5-VL-3B-Instruct as the baseline, reproduced it
on TextVQA val (79.60 against the published 79.3), measured the zero-shot baseline on ViVQA
and ViTextVQA, and then adapted it with QLoRA using 7.37 million trainable parameters, i.e.\
0.196\% of all parameters.

The proposed improvement is OCR-aware QLoRA: the text in the image is read with VietOCR,
sorted in reading order and added to the prompt already during training. On the full
ViTextVQA test set of 10,028 questions, the method reaches 53.35 EM, 3.09 points above plain
QLoRA on the same data ($p = 6.7\times10^{-21}$) and above the best previously published
result (25.48 EM). On ViVQA, QLoRA raises EM from 22.09 to 76.64.

%********************************** Section 6.2 ********************************
\section{Answers to the research questions}
\label{sec:con-rq}

\begin{description}
  \item[RQ1.] \textit{How well does zero-shot perform, and which Vietnamese errors dominate?}
  Zero-shot reaches 22.09 EM on ViVQA and 29.12 on ViTextVQA. The dominant error is the wrong
  format: the model answers at greater length than the gold answer (2.9 against 1.6 words on
  ViVQA) and sometimes in English. Answers with tone marks are clearly harder than answers
  made only of Latin letters. A few-shot formatting prompt does not fix these errors (H3 not
  supported).

  \item[RQ2.] \textit{How much does QLoRA improve without full fine-tuning?}
  QLoRA increases EM by 21.1 points on ViTextVQA and 54.6 points on ViVQA (the official model
  C3), with only 0.196\% of the parameters, one L4 GPU, at most 12.75 hours of training and a
  peak VRAM below 8.4~GB. The gain comes mainly from learning the right answer format and is
  not due to image leakage: on the clean test set, the gain is still 44.1 points (H1
  supported).

  \item[RQ3.] \textit{Does OCR text help, and does it work together with QLoRA?}
  Yes, but only when the model is trained with OCR. Adding OCR to the prompt of a model that
  has never seen OCR brings no improvement; training with OCR gives an interaction effect of
  $+4.9$ EM points, stable across three seeds, and the benefit does not shrink with more data.
  The ANLS gain ($+1.87$) is below the pre-set 3-point threshold, so H2 is partially
  supported: OCR helps produce exactly the right string more than it helps produce
  character-level near-misses.

  \item[RQ4.] \textit{After adaptation, does the model invent content more or less often?}
  On POPE-vi, B3-full does not increase hallucination according to the pre-set rule
  ($\Delta$F1 from $+2.5$ to $+5.5$), but its object false-alarm rate increases under
  adversarial sampling. C3 cannot be measured because it has forgotten how to answer yes/no.
  The manual review shows that hallucination accounts for only about 2.3\% (B3-full) and 3.7\%
  (C3) of all answers. H4 is partially supported.
\end{description}

%********************************** Section 6.3 ********************************
\section{Contributions}
\label{sec:con-contrib}

\begin{enumerate}
  \item \textbf{OCR-aware QLoRA for Vietnamese VQA}, with its mechanism isolated by a
  $2\times2$ experiment: the benefit comes from teaching the model to use OCR, not from adding
  OCR to the prompt.
  \item \textbf{Efficient and reproducible adaptation.} The results on ViTextVQA are above the
  previously published results with 0.196\% of the parameters on one GPU; every number has a
  log, a seed and a configuration, with a reproducibility package containing the source code,
  adapters, configurations, inference scripts and instructions to re-run.
  \item \textbf{Measuring problems specific to Vietnamese.} Each of the four problems V1--V4
  has concrete measurements; POPE-vi is a Vietnamese hallucination benchmark of 4,985
  questions on 500 images unseen in training.
  \item \textbf{Two-level error analysis.} An automatic analysis by error type and OCR status,
  checked against a blind manual review, shows that EM underestimates the model and estimates
  how much could still be gained with better OCR (about 17 EM points).
  \item \textbf{A verified prototype.} A web application that answers with the OCR text and a
  confidence level, verified to give the same answers as the official evaluation pipeline,
  with a median latency of about 1 second on one L4 GPU.
\end{enumerate}

%********************************** Section 6.4 ********************************
\section{Limitations}
\label{sec:con-limits}

The main limitations were analysed in Section~\ref{sec:dis-validity}: the main models were run
with only one seed; EM is too strict towards answers that are correct in meaning but written
differently; the quality of Vietnamese scene-text OCR is a bottleneck; the ANLS threshold of H2
was not reached; F1 on POPE-vi hides the false-alarm rate; the comparison with published results
depends on the paper's answer normalisation; and the final manual review had only one
annotator.

%********************************** Section 6.5 ********************************
\section{Future work}
\label{sec:con-future}

\begin{enumerate}
  \item \textbf{Better Vietnamese scene-text OCR}: fine-tune VietOCR on scene text and improve
  the reading order. According to the error analysis, this direction has the highest ceiling
  for improvement.
  \item \textbf{Position-aware OCR}: feed the coordinates and order of the text boxes to the
  model in the spirit of M4C~\cite{hu2020m4c}, to reduce wrong phrase boundaries.
  \item \textbf{Human evaluation and a synonym-tolerant metric}, since more than half of the EM
  errors are correct in meaning.
  \item \textbf{Balancing yes/no questions during training} to avoid forgetting the format as
  in C3, and measuring text hallucination on images with text, not only object hallucination.
  \item \textbf{Confidence calibration} (e.g.\ temperature scaling) so that the prototype can
  show probabilities instead of levels.
  \item \textbf{More seeds for the main models} and newer base models under the same evaluation
  protocol.
\end{enumerate}

The results of this thesis show that a team with limited resources can still adapt a large VLM
to Vietnamese and surpass previous work, provided that every change is measured under the same
conditions and every conclusion is tested.
